\Hefteintrag{1.7}{Programmieren: Sprachen}{
Beim Programmieren schreibt man präzise und unmissverständliche Befehle, die dann vom \LoesungLuecke{Compiler}{6cm}~in Maschinensprache übersetzt werden. 

Hierfür gibt es verschiedene Programmiersprachen mit jeweils einer festgelegte \emphColA{Grammatik und Rechtschreibung (=\LoesungLuecke{Syntax}{4cm})}. Dem gegenüber steht die \emphColC{Logik eines Programms (=~\LoesungLuecke{Semantik}{4cm})}, die weitgehend unabhängig von der Programmiersprache ist.

\vspace{0.2cm}
Bei den einsteigerfreundlichen \emphColA{block-basierten Sprachen wie Blockly oder Online-Karol} übernehmen die Blöcke und ihre Verbindungsmöglichkeiten weitgehend die Einhaltung des Syntax. Bei \emphColC{text-basierten Sprachen wie \LoesungLuecke{Java}{6cm}~oder \LoesungLuecke{Python}{6cm}} müssen die Entwickler:innen selbst auf deren Einhaltung achten. Diese Sprachen bieten dadurch aber auch weitaus mehr Möglichkeiten (\emph{man sagt auch: sie sind mächtiger}). 

Dieses Skript arbeitet parallel mit den blockbasierten Sprachen Blockly und Online-Karol und der weit verbreiteten text-basierten und \emphColA{objektorientierten} Sprache \emphColA{Java}.

}